\section{Acesso ao Repositório}

\subsection{Localização}

O repositório do projeto encontra-se disponível publicamente no GitHub:

\begin{center}
\url{https://github.com/MAlves4018/NatureProtector.git}
\end{center}

Por se tratar de um repositório público, a consulta do código fonte, documentação e artefactos versionados não requer permissões adicionais de leitura.

\subsection{Clonagem do repositório}

A obtenção local do repositório pode ser feita com:

\begin{verbatim}
git clone https://github.com/MAlves4018/NatureProtector.git
cd NatureProtector
\end{verbatim}

Após a clonagem, recomenda-se confirmar a branch e o commit em uso:

\begin{verbatim}
git branch --show-current
git log -1 --oneline
git status --short
\end{verbatim}

Estas informações permitem associar qualquer evidência de execução, teste ou demonstração a uma versão concreta do repositório.

\subsection{Setup local}

O setup local não é duplicado neste documento. A referência de setup é:

\begin{center}
\repoPath{docs/local-baseline-setup.md}
\end{center}

Esse ficheiro deve ser lido antes de tentar executar a baseline local. Em termos gerais, o repositório depende de um ambiente com .NET SDK, Docker Desktop, PowerShell e Node.js/npm para executar a solução, a infraestrutura local e a interface web.
